\scanpage{174}
% archive foliation: 173
Il en résulte en effet que l'on a une suite exacte (en se bornant aux faisceaux inversibles)
\diagram{\begin{tikzcd}[column sep=small]
0\arrow[r]&\operatorname{Pic}(S)\arrow[r]&\operatorname{Pic}(X)\arrow[r]&\mathrm{Prépic}'(X/S)
\end{tikzcd}}
l'image \inserted{dernière} étant formée des sections de $\underline H^1_{X/S}(\underline O_X^*)$ qui sont dans l'image de $\underline H^1(\underline O_X^*)$ i.e. dont l'obstruction dans $H^2(S,\underline O_S^*)$ est nulle.
\margin{[lorsque $X/S$ a une section] \unsure{les hom sont définis en tous cas}, mais l'exactitude demande que $\underline H^0_{X/S}(\underline O_X)=\underline O_S$}

[N.B. le début de la suite exacte n'est que la suite exacte de bord de la suite spectrale de Leray
\diagram{\begin{tikzcd}[column sep=small]
0\arrow[r]&H^1(S,\underline H^0(\underline O_X^*))\arrow[r]&H^1(X,\underline O_X^*)\arrow[r]&H^0(S,\underline H^1(\underline O_X^*))\arrow[r]&H^2(S,\underline H^0(\underline O_X^*))\arrow[d]\\
&&&&H^2(X,\underline O_X^*)
\end{tikzcd}}
compte tenu que $\underline H^0(\underline O_X^*)=\underline O_S^*$]. S'il y a une section $\inserted{s}$ de $X/S$, alors on \deleted{\unsure{aura un bien inverse}} a un résultat plus complet

\uline{Corollaire 2} Soit $\underline C$ la catégorie des faisceaux inversibles $\underline L$ sur $X$ munis d'un \uline{isomorphisme} $\underline O_S\to s^*(\underline L)$. Alors \deleted{pour}\inserted{(i)} tout automorphisme d'un objet de $\underline C$ est trivial.

(ii) Pour \deleted{tout} tout élém$^t$ de $\mathrm{Prépic}(X/S)$ provient d'un $\underline L\in\underline C$, déterminé à isom. unique près.

